% ─── Capítulo 2: Analizador Sintáctico ────────────────────────────────────────
\chapter{FRONTEND — Análisis Sintáctico}

\section{Fundamentos Teóricos}

El \textbf{análisis sintáctico} (o \textit{parsing}) es la segunda etapa de la compilación. A diferencia del lexer, que agrupa caracteres en tokens planos, el parser descubre la \textbf{estructura jerárquica} del programa. Comprueba si la secuencia lineal de tokens obedece a las reglas estructurales del lenguaje.

Teóricamente, el lenguaje se modela mediante una \textbf{Gramática Libre de Contexto (CFG)}, definida por:
\begin{itemize}
    \item \textbf{Terminales:} Los tokens provenientes del lexer.
    \item \textbf{No Terminales:} Variables sintácticas (ej. \textit{expresion}, \textit{declaracion}).
    \item \textbf{Producciones:} Reglas de reemplazo (\textit{No Terminal $\rightarrow$ símbolos}).
    \item \textbf{Símbolo Inicial:} El punto de partida (ej. \textit{programa}).
\end{itemize}

\begin{figure}[H]
    \centering
    \begin{tikzpicture}[
        level distance=1.2cm,
        level 1/.style={sibling distance=4cm},
        level 2/.style={sibling distance=2cm},
        level 3/.style={sibling distance=1.5cm},
        every node/.style={circle, draw, fill=tppblue!10, text=tppdark, minimum size=0.8cm, font=\small, thick},
        edge from parent/.style={thick, draw=tppblue}
        ]
        \node {expr}
            child {node {expr}
                child {node[fill=tppgreen!10] {NUM}}
            }
            child {node[fill=tppgray!30, draw=gray] {+}}
            child {node {expr}
                child {node {expr}
                    child {node[fill=tppgreen!10] {ID}}
                }
                child {node[fill=tppgray!30, draw=gray] {*}}
                child {node {expr}
                    child {node[fill=tppgreen!10] {NUM}}
                }
            };
    \end{tikzpicture}
    \caption{Ejemplo de Árbol de Derivación (Parse Tree) para la expresión matemática $3 + x * 5$, resolviendo la precedencia mediante la gramática.}
    \label{fig:parse_tree}
\end{figure}

\begin{infobox}[¿Por qué es necesario el Análisis Sintáctico?]
Una gramática regular (usada en el lexer) no puede contar ni anidar estructuras (no puede verificar que los paréntesis o llaves estén correctamente balanceados). Una CFG sí puede hacerlo. El análisis sintáctico permite:
\begin{enumerate}
    \item Validar que la sintaxis general es correcta (ej. un \texttt{si} seguido de \texttt{()} y \texttt{\{\}}).
    \item Capturar precedencia y asociatividad de operadores.
    \item Construir una representación en memoria abstracta (AST) para manipular el programa más fácilmente en el futuro.
\end{enumerate}
\end{infobox}

\section{Implementación con Bison: \texttt{parser.y}}

Para construir el parser de TPP se emplea \textbf{Bison}, que es un generador de parsers de tipo \textbf{LALR(1)} (\textit{Look-Ahead Left-to-Right, Rightmost derivation with 1 token lookahead}). Bison toma la CFG escrita en \texttt{parser.y} y genera una máquina de estados (Pushdown Automaton) en C (\texttt{parser.tab.c}).

\subsection{Funcionamiento Interno (Shift/Reduce)}

El parser generado por Bison funciona mediante dos operaciones principales en una pila:
\begin{itemize}
    \item \textbf{Shift (Desplazamiento):} Lee un token de la entrada y lo empuja a la pila sintáctica, pasando a un nuevo estado.
    \item \textbf{Reduce (Reducción):} Cuando la cima de la pila coincide con el lado derecho de una producción, saca esos elementos de la pila y empuja el lado izquierdo (el No Terminal). En este momento se ejecuta el código C asociado (la acción semántica), que en nuestro caso \textbf{crea un nodo del AST}.
\end{itemize}

\subsection{Gramática del Lenguaje TPP}

La gramática completa se define mediante reglas BNF en el archivo \texttt{parser.y}:

\begin{lstlisting}[style=cstyle, caption={Reglas gramaticales principales}]
// Punto de entrada
programa:
    lista_elementos { 
        NodoAST *raiz = nuevo_nodo_programa($1);
        // Disparar el pipeline completo de compilación
        analizar_semantica(raiz);
        optimizar_ast(raiz);
        generar_codigo_intermedio(raiz);
        generar_codigo(raiz, out);
    }

// Declaración de función
declaracion_fun:
    FUN ID PARI parametros PARD FLECHA tipo LLAVEI 
    bloque_instrucciones LLAVED
    { $$ = nuevo_nodo_declaracion_fun($2, $7, $4, $9); }
    
    | FUN ID PARI parametros PARD LLAVEI 
    bloque_instrucciones LLAVED
    { $$ = nuevo_nodo_declaracion_fun($2, 0, $4, $7); }

// Declaracion de variable con tipo
declaracion_var:
    tipo lista_ids PUNTOCOMA
    { $$ = nuevo_nodo_declaracion_var($1, $2); }

// Estructura si/sino
estructura_si:
    SI PARI expresion PARD LLAVEI bloque_instrucciones LLAVED
    { $$ = nuevo_nodo_si($3, $6, NULL); } %prec LOWER_THAN_SINO
    
    | SI PARI expresion PARD LLAVEI bloque_instrucciones LLAVED
      SINO LLAVEI bloque_instrucciones LLAVED
    { $$ = nuevo_nodo_si($3, $6, $10); }

// Llamada a función
llamada_funcion:
    ID PARI argumentos PARD 
    { $$ = nuevo_nodo_llamada_funcion($1, $3); }

// Expresiones con precedencia
expresion:
    expresion MAS expresion  { $$ = nuevo_nodo_operacion(MAS, $1, $3); }
    | expresion MENOS expresion { $$ = nuevo_nodo_operacion(MENOS, $1, $3); }
    | expresion MULT expresion  { $$ = nuevo_nodo_operacion(MULT, $1, $3); }
    | expresion DIV expresion   { $$ = nuevo_nodo_operacion(DIV, $1, $3); }
    | expresion IGUAL expresion { $$ = nuevo_nodo_operacion(IGUAL, $1, $3); }
    | expresion MENOR expresion { $$ = nuevo_nodo_operacion(MENOR, $1, $3); }
    | NUM_ENTERO  { $$ = nuevo_nodo_entero($1); }
    | NUM_DECIMAL { $$ = nuevo_nodo_decimal($1); }
    | ID          { $$ = nuevo_nodo_variable($1); }
    | CADENA      { $$ = nuevo_nodo_cadena($1); }
    | llamada_funcion { $$ = $1; }
\end{lstlisting}

\subsection{Resolución de Ambigüedades}

Bison resuelve ambigüedades mediante reglas de precedencia declaradas explícitamente:

\begin{lstlisting}[style=cstyle, caption={Declaraciones de precedencia}]
// Menor a mayor precedencia (de arriba hacia abajo)
%left IGUAL DIFERENTE MENOR MAYOR MENOR_IGUAL MAYOR_IGUAL
%left MAS MENOS
%left MULT DIV MOD

// Resolución del "dangling else" (sino huérfano)
%nonassoc LOWER_THAN_SINO
%nonassoc SINO
\end{lstlisting}

Esto garantiza que \texttt{3 + 4 * 2} se evalúe como \texttt{3 + (4 * 2) = 11} y no como \texttt{(3 + 4) * 2 = 14}.

\section{Estructuras Sintácticas Soportadas}

\subsection{Declaraciones de Variables}

\begin{lstlisting}[style=tppstyle, caption={Formas válidas de declaración}]
entero x;                 // Simple sin inicializacion
entero x = 5;             // Con inicializacion
entero x = 3 + 2;         // Con expresion
entero arr[10];           // Arreglo de 10 enteros
decimal pi = 3.14;        // Variable decimal
\end{lstlisting}

\subsection{Declaraciones de Funciones}

\begin{lstlisting}[style=tppstyle, caption={Formas de declaración de funciones}]
// Función sin retorno (void)
fun saludar() {
    imprimir("Hola");
}

// Función con parámetros y tipo de retorno
fun sumar(entero a, entero b) -> entero {
    retornar a + b;
}

// Función main (punto de entrada del usuario)
fun main() {
    entero resultado;
    resultado = sumar(3, 4);
    imprimir(resultado);
}
\end{lstlisting}

\subsection{Estructuras de Control}

\begin{lstlisting}[style=tppstyle, caption={Estructuras de control disponibles}]
// Condicional
si (x > 0) {
    imprimir("positivo");
} sino {
    imprimir("negativo o cero");
}

// Bucle mientras
mientras (i < 10) {
    i = i + 1;
}

// Bucle para (con declaracion de variable en cabecera)
para(entero i = 0; i < 5; i = i + 1) {
    imprimir(arr[i]);
}

// Selección múltiple
depende (x) {
    caso 1: { imprimir("uno"); }
    caso 2: { imprimir("dos"); }
    otros:  { imprimir("otro"); }
}
\end{lstlisting}

\section{Manejo de Errores Sintácticos}

El parser implementa \textbf{recuperación de errores} para continuar el análisis aunque encuentre errores:

\begin{lstlisting}[style=cstyle, caption={Recuperación de errores en el parser}]
estructura_si:
    SI PARI expresion PARD LLAVEI bloque_instrucciones LLAVED { ... }
    | SI PARI error PARD LLAVEI bloque_instrucciones LLAVED {
        yyerrok;
        printf("=> [Panic Mode] Error en condicion del SI.\n");
        $$ = NULL;
    }

estructura_para:
    PARA PARI tipo ID ASIG expresion PUNTOCOMA 
    expresion PUNTOCOMA ID ASIG expresion PARD 
    LLAVEI bloque_instrucciones LLAVED { ... }
    | PARA PARI error PARD LLAVEI bloque_instrucciones LLAVED {
        yyerrok;
        $$ = NULL;
    }
\end{lstlisting}

La función \texttt{yyerror()} produce mensajes descriptivos con línea y columna:

\begin{verbatim}
¡Ups! Hay un problema sintáctico en tu código.
Linea 14, Columna 11
Detalle: syntax error cerca del elemento 'cinco'
\end{verbatim}

\section{Orquestación del Pipeline}

Una vez que el parser construye el AST sin errores, el punto de entrada \texttt{programa} actúa como orquestador y dispara todas las fases sucesivas del compilador en secuencia:

\begin{lstlisting}[style=cstyle, caption={Orquestación del pipeline en parser.y}]
programa:
    lista_elementos {
        NodoAST *raiz = nuevo_nodo_programa($1);
        
        // Fase 1: Análisis Semántico
        int errores_sem = analizar_semantica(raiz);
        
        if (errores_sem == 1) { // Éxito
            // Fase 2: Optimización del AST
            optimizar_ast(raiz);
            
            // Fase 3: Código Intermedio
            generar_codigo_intermedio(raiz);
            
            // Fase 4: Generación de Código Final
            FILE *out = fopen("output.s", "w");
            generar_codigo(raiz, out);
            fclose(out);
        }
    }
\end{lstlisting}
